\chapter{Wavefront aberrations and what RMS does not tell you}\label{ch:aberration}
\paperloc{Introduction, Section 2, and the motivation for an imaging-based budget.}
\reading{Born--Wolf \S5.1, printed pp.~203--206 (begin at PDF p.~232);
\S9.1.2, p.~462 (PDF p.~557); \S9.1.3, pp.~463--464
(PDF pp.~556--555); \S9.3, pp.~468--469 (PDF pp.~551--550).
The backward PDF order is real.}

\section{The reference sphere and the aberration function}
An ideal converging wave reaches its intended image point with the same
phase contribution from every admitted pupil location. On a suitable
reference sphere centered at that image point, its phase is constant.
The actual wave has residual phase variation $2\pi W/\lambda$.
The real-valued function $W(\rho,\theta)$ is the wave aberration.

The coordinate $(\rho,\theta)$ identifies a point in the pupil, not
a point on the wafer. A pupil map is therefore not a picture of an
aberrated contact hole. To get that picture, multiply the pupil field by
the phase factor and perform the diffraction calculation.

At a different field point on the wafer, the same optical system may
have a different pupil wavefront:
\[
W=W(\rho,\theta;\vct H),
\]
where $\vct H$ denotes image-field position. A set of coefficients for
one field point is not automatically valid across an entire exposure
field. Likewise, scan time and thermal state may add dependencies.

\section{The familiar named aberrations}
For a centred rotationally symmetric system, a low-order expansion
at field height $H$ contains terms schematically of the form
\begin{align}
W={}&A\rho^4
+B H\rho^3\cos\theta
+C H^2\rho^2\cos^2\theta\\
&+D H^2\rho^2
+E H^3\rho\cos\theta+\cdots.
\label{eq:seidel}
\end{align}
The coefficients and axes depend on convention, but the shapes reveal
the standard families:
\begin{itemize}
\item \textbf{Spherical aberration:} different pupil radii have different
focal behavior; the wavefront has a quartic radial contribution.
\item \textbf{Coma:} an off-axis asymmetric wavefront component, often
producing a lopsided point image.
\item \textbf{Astigmatism:} different meridional directions focus differently.
At a suitable intermediate focus, its wavefront has a
$\rho^2\cos2\theta$ or $\rho^2\sin2\theta$ component.
\item \textbf{Field curvature:} best focus varies across the image field.
\item \textbf{Distortion:} field-dependent placement error; a point may
remain locally sharp while its position is incorrect.
\end{itemize}
Piston, tilt, and defocus specify a constant phase, a linear pupil phase,
and a quadratic radial phase. Whether they are called aberrations or
alignment freedoms depends on the measurement purpose. The paper includes
low-order coefficients because image placement and focus behavior matter.

Equation \eqref{eq:seidel} is an illustrative classical expansion, not the
prescription of the paper's off-axis multi-mirror projector. A real EUV
system need not have the centred-system symmetries used to derive this form.

\section{Piston does not change intensity}
If $W$ is replaced by $W+c_0$, then every field contribution receives the
same factor:
\[
U_{\rm image}\mapsto\ee^{\ii2\pi c_0/\lambda}U_{\rm image}.
\]
The intensity is unchanged because
$|\ee^{\ii2\pi c_0/\lambda}|^2=1$.
Consequently intensity-derived CD, PS, and PVB cannot recover global
piston. This is an exact invariance within the model, not a failure of
neural-network training.

\section{Tilt moves the image}
Consider a linear OPD in the normalized frequency pupil,
$W=t_xu+t_yv$. Then
\[
P(\vct f)=P_0(\vct f)
\exp\!\left[\ii2\pi\left(\frac{t_x}{\NA}f_x
+\frac{t_y}{\NA}f_y\right)\right].
\]
Substituting in the inverse Fourier integral shows
\begin{equation}
U(\vct x)=U_0\!\left(\vct x+\frac{\vct t}{\NA}\right),
\qquad
\Delta\vct x=-\frac{\vct t}{\NA}.
\label{eq:tilt}
\end{equation}
Thus ideal pure tilt shifts a complete image without changing its shape.
An image registered back to its own centre has the same CD and contrast.
An unregistered image has a nonzero placement error.

The sign follows our transfer-function convention. The measurable
magnitude and the distinction between shape and placement are the
important points. When metrology removes tilt before quoting RMS,
that RMS cannot subsequently bound the removed placement error.

\section{Defocus changes the reference image plane}
Equation \eqref{eq:defocusphase} gives a quadratic pupil phase when the
observation plane moves. Equivalently, a quadratic aberration can be partly
compensated by changing focus. A report of best-focus RMS may therefore
exclude a defocus term that still matters at a fixed wafer plane.

The paper varies the focus setting by $\pm15\nm$ while also changing
aberration coefficients. To reconstruct this experiment one must know
whether the Zernike coefficients are defined at a fixed reference focus,
whether refocusing is permitted, and whether focus is reapplied separately
for every candidate. The publication does not fully specify this convention.

\section{Deriving the small-aberration Strehl relation}
For a clear, uniformly illuminated pupil and an on-axis point source,
the normalized field at the chosen reference image point is
\[
\frac{U(0)}{U_{\rm ideal}(0)}
=\avg{\ee^{\ii\phi}},\qquad \phi=\frac{2\pi W}{\lambda}.
\]
Expand the exponential:
\[
\avg{\ee^{\ii\phi}}\simeq
1+\ii\avg{\phi}-\frac12\avg{\phi^2}.
\]
Multiplying by its conjugate and retaining terms through second order,
\begin{equation}
S_{\rm ref}\simeq1-\avg{\phi^2}+\avg{\phi}^2
=1-\left(\frac{2\pi\sigma_W}{\lambda}\right)^2.
\label{eq:strehl}
\end{equation}
The exponential form often used for small wavefront errors is
\begin{equation}
S\simeq\exp[-(2\pi\sigma_W/\lambda)^2].
\label{eq:marechal}
\end{equation}
It shares the leading expansion. It is also the exact squared modulus
of the characteristic function for a zero-mean Gaussian phase distribution
under the corresponding average. An arbitrary deterministic aberration
does not make that exponential exact.

The usual Strehl ratio compares peak intensities under consistent total
power and aperture conditions. Equation \eqref{eq:strehl} initially evaluates
a chosen reference point. To discuss the optimized peak, remove or optimize
the appropriate tilt and focus and specify that choice.
With a nonuniform pupil, the amplitude weighting in the coherent field
average must also be treated explicitly.

\worked{Interpreting a $\lambda/50$ RMS requirement}
At 13.5 nm, $\lambda/50=0.27\nm=20\,\mathrm{m}\lambda$.
If this is piston-removed RMS in a suitable small-aberration setting,
\[
\sigma_\phi=2\pi/50=0.125664,\qquad
S\simeq\ee^{-(2\pi/50)^2}=0.98433.
\]
A high Strehl ratio describes concentration of a point image. It does
not guarantee a subnanometre PS for every periodic contact image.
Pure tilt illustrates the distinction particularly clearly.

\section{Equal RMS is not equal lithographic performance}
RMS measures a mean-square phase error over the pupil. A periodic mask
uses particular sets of pupil locations and their phase differences.
Two maps with identical RMS can give different phase differences at
those locations. Their CD, PS, and process window can therefore differ.

This is the physical motivation for the paper. It does not imply that
RMS is useless. RMS remains a compact measure of optical quality and a
useful control variable. It means that RMS alone is not a complete
description of a pattern-specific imaging requirement.

\chapter{Zernike polynomials from the beginning}\label{ch:zernike}
\paperloc{Equation (1), Table 1, and the coefficient axes in Figs.~5, 9, 11, and 12.}
\reading{Born--Wolf \S9.2.1, printed pp.~464--466
(PDF pp.~555, 554, 553), including its radial-polynomial table.
Its normalization differs from the RMS-normalized real basis used here.}

\section{Why use orthogonal functions?}
Suppose we tried to describe $W$ with $1,\rho^2,\rho^4,\ldots$.
These functions overlap under disk averaging:
$\avg{\rho^2}=1/2$ and $\avg{\rho^2\rho^4}=1/4$, not zero.
Changing one fitted coefficient can compensate another.
Orthogonal functions separate those overlapping contributions.

Define the inner product
\[
(f,g)=\avg{fg}
=\frac1\pi\int_0^{2\pi}\int_0^1fg\,\rho\dd\rho\dd\theta.
\]
Orthogonal means $(f,g)=0$ for different modes.
Orthonormal means each mode also has $(f,f)=1$.
The analogy is perpendicular unit axes in ordinary Euclidean geometry,
but the vectors here are functions over a disk.

\section{Constructing the first modes explicitly}
The constant mode is $Z_0^0=1$. Its RMS is one.
For a tilt, $\rho\cos\theta=u$ and $\avg{u^2}=1/4$,
so the unit-RMS tilt is
\[
Z_1^1=2\rho\cos\theta,\qquad
Z_1^{-1}=2\rho\sin\theta.
\]
To create defocus orthogonal to piston, subtract the mean from $\rho^2$:
$\rho^2-1/2$. Its mean square is
$1/3-1/4=1/12$, giving
\begin{equation}
Z_2^0=\sqrt3(2\rho^2-1).
\end{equation}
For astigmatism, angular averaging already removes piston:
\[
\avg{(\rho^2\cos2\theta)^2}
=\frac12\avg{\rho^4}=\frac16.
\]
Hence
$Z_2^2=\sqrt6\rho^2\cos2\theta$ and
$Z_2^{-2}=\sqrt6\rho^2\sin2\theta$.

For a coma-shaped mode, start from
$(\rho^3-a\rho)\cos\theta$. Remove its overlap with tilt:
\[
0=\avg{(\rho^3-a\rho)\rho\cos^2\theta}
=\frac12\left(\frac13-\frac a2\right),
\]
which gives $a=2/3$. Rescale the radial polynomial to equal one at
the rim, giving $3\rho^3-2\rho$. Its normalized real mode is
\begin{equation}
Z_3^1=\sqrt8(3\rho^3-2\rho)\cos\theta.
\label{eq:coma}
\end{equation}
The sine partner is $Z_3^{-1}$. Notice the subtracted tilt. A balanced
Zernike coma mode is not just the monomial $\rho^3\cos\theta$.

For spherical aberration, start with $\rho^4+a\rho^2+b$ and require
orthogonality to $1$ and $\rho^2$:
\[
\frac13+\frac a2+b=0,\qquad
\frac14+\frac a3+\frac b2=0.
\]
Solving gives $a=-1$, $b=1/6$. Rim normalization and unit RMS produce
\begin{equation}
Z_4^0=\sqrt5(6\rho^4-6\rho^2+1).
\label{eq:spherical}
\end{equation}
This explains why the spherical Zernike mode contains defocus and piston
terms: they make it independent of the lower radial modes.

\section{The general definition}
Let $n\geq0$ be radial degree and $m$ be an integer with
$|m|\leq n$ and $n-|m|$ even. Define
\begin{equation}
R_n^{|m|}(\rho)=
\sum_{s=0}^{(n-|m|)/2}
(-1)^s
\frac{(n-s)!}
{s!\,[(n+|m|)/2-s]!\,[(n-|m|)/2-s]!}
\rho^{n-2s}.
\label{eq:radial}
\end{equation}
The allowed parity ensures all factorials have nonnegative integer
arguments and all powers have the correct angular symmetry.
The rim value is $R_n^{|m|}(1)=1$.

Our real RMS-normalized convention is
\begin{equation}
Z_n^m=
\begin{cases}
\sqrt{n+1}\,R_n^0,&m=0,\\
\sqrt{2(n+1)}\,R_n^m\cos(m\theta),&m>0,\\
\sqrt{2(n+1)}\,R_n^{|m|}\sin(|m|\theta),&m<0.
\end{cases}
\label{eq:Zdef}
\end{equation}
The radial orthogonality identity is
\[
\int_0^1R_n^mR_{n'}^m\,\rho\dd\rho
=\frac{\delta_{nn'}}{2(n+1)}.
\]
Angular integration gives $2\pi$ for constant modes and $\pi$ for
the square of each nonconstant sine or cosine. Combining those factors
with Eq.~\eqref{eq:Zdef} gives
\begin{equation}
\avg{Z_n^mZ_{n'}^{m'}}=\delta_{nn'}\delta_{mm'}.
\label{eq:Zorth}
\end{equation}
The Kronecker delta $\delta_{ab}$ equals one if $a=b$ and zero otherwise.
This is the formal statement that the modes form independent unit axes
under this pupil-area metric.

\section{A radial table sufficient to understand low-order budgets}
Every $m>0$ entry below has both a cosine and sine partner. Apply the
normalization in Eq.~\eqref{eq:Zdef}; the table lists only radial factors.
\begin{longtable}{@{}ccp{.64\textwidth}@{}}
\toprule $n$ & $|m|$ & $R_n^{|m|}(\rho)$\\ \midrule\endhead
0&0&$1$\\
1&1&$\rho$\\
2&0&$2\rho^2-1$\\
2&2&$\rho^2$\\
3&1&$3\rho^3-2\rho$\\
3&3&$\rho^3$\\
4&0&$6\rho^4-6\rho^2+1$\\
4&2&$4\rho^4-3\rho^2$\\
4&4&$\rho^4$\\
5&1&$10\rho^5-12\rho^3+3\rho$\\
5&3&$5\rho^5-4\rho^3$\\
5&5&$\rho^5$\\
6&0&$20\rho^6-30\rho^4+12\rho^2-1$\\
6&2&$15\rho^6-20\rho^4+6\rho^2$\\
6&4&$6\rho^6-5\rho^4$\\
6&6&$\rho^6$\\
\bottomrule
\end{longtable}
There are $n+1$ real modes at degree $n$. Through degree six there are
$1+2+\cdots+7=28$ modes, including piston.
This is a complete sequence by radial degree.
It is deliberately not presented as the paper's unspecified single-index
list of 25 terms.

\figteach{zernike_atlas}{.93}{Original atlas of selected real,
RMS-normalized modes. Red and blue indicate opposite OPD signs; the same
colour scale is used throughout. The labels are $(n,m)$, not the paper's
single-index coefficient numbers.}

\section{Recovering coefficients and calculating RMS}
Expand a square-integrable wavefront as
\begin{equation}
W(\rho,\theta)=\sum_{n,m}c_n^mZ_n^m(\rho,\theta).
\end{equation}
Multiply by $Z_{n'}^{m'}$ and average. Orthogonality removes all other
terms:
\begin{equation}
c_{n'}^{m'}=\avg{WZ_{n'}^{m'}}.
\label{eq:projection}
\end{equation}
Each coefficient has the units of $W$. If $W$ is in nanometres, so is $c$.
If $W/\lambda$ is expanded, the coefficients are in waves.

Because $\avg{Z_0^0}=1$ and all other mode means are zero,
$\avg W=c_0^0$. Expanding the square and using orthogonality gives
\begin{equation}
\sigma_W^2=\sum_{(n,m)\ne(0,0)}(c_n^m)^2.
\label{eq:rss}
\end{equation}
This identity depends on the unit-RMS normalization, complete disk,
uniform area weighting, and explicit removal of piston.
Other removed modes must also be excluded. With an unnormalized
basis, the corresponding modal norm factors remain in the sum.

\worked{A coefficient is not a peak-to-valley value}
For defocus $W=cZ_2^0$, the centre value is $-\sqrt3c$ and the rim value
is $+\sqrt3c$. Thus
$\mathrm{PV}=2\sqrt3|c|$, while $\sigma_W=|c|$.
For $c=0.27\nm$, RMS is 0.27 nm but PV is about 0.9353 nm.
For the normalized tilt $W=2cu$, PV is $4|c|$.
Equal coefficients mean equal RMS for these modes, not equal PV.

\worked{Why per-mode limits are not a total RMS limit}
If 24 independent orthonormal, nonpiston modes each have magnitude
$20\,\mathrm{m}\lambda$, then
\[
\sigma_W=20\sqrt{24}\,\mathrm{m}\lambda
=97.9796\,\mathrm{m}\lambda=1.32272\nm.
\]
The input box $|c_j|\leq20\,\mathrm{m}\lambda$ therefore does not mean
total RMS is at most $\lambda/50$. If all coefficients are independent
uniform variables on that interval, then
$\E[c_j^2]=(20\,\mathrm{m}\lambda)^2/3$, and
\[
\sqrt{\E[\sigma_W^2]}=20\sqrt{24/3}
=56.5685\,\mathrm{m}\lambda.
\]
This is the square root of the expected squared RMS, not the expectation
of RMS itself. Both calculations are conditional teaching results:
the paper does not state the normalization needed to apply them directly
to its coefficient arrays.

\section{Single-index numbering is a convention, not a law}
Many programs replace $(n,m)$ by a single integer $j$.
Noll, Fringe, OSA/ANSI, and software-specific sequences differ.
Some start counting at zero, some at one; some normalize each mode to
unit RMS and some do not. Axes, sine/cosine assignment, and sign can differ.

The paper writes
\[
W(\rho,\theta)=\sum_i Z_iF_i(\rho,\theta).
\]
Here its $Z_i$ are \emph{coefficients}, while its $F_i$ are basis
functions. In this book $Z_n^m$ names a basis function and $c_n^m$ its
coefficient. These are two notation systems for the same expansion idea.
Confusing them changes the meaning of an equation.

The paper does not supply the explicit $F_i$ table or normalization.
We therefore cannot safely identify its $Z_{20}$, $Z_{21}$, or $Z_{22}$
with particular named modes. Its report that these coefficients have
narrow ranges is meaningful in its own simulator coordinates, but a
physical mode-name interpretation requires that missing dictionary.

\section{Mode orientation and rotating axes}
For a sine/cosine pair of the same $n$ and $m>0$,
\[
c_c\cos m\theta+c_s\sin m\theta
=C\cos[m(\theta-\theta_0)],
\]
where
$C=\sqrt{c_c^2+c_s^2}$ and
$\theta_0=\operatorname{atan2}(c_s,c_c)/m$.
The coefficient pair rotates when axes rotate, but $C$ stays unchanged.
Astigmatism has $m=2$, so a $45^\circ$ rotation in physical axes
interchanges the character of its cosine and sine partners.
This matters when comparing horizontal and vertical pattern sensitivities.

\section{Balanced spherical aberration as a minimization}
Suppose a surface produces raw quartic OPD $B\rho^4$. Permit a compensating
defocus $d\rho^2$ and remove piston afterward. Its variance is
\[
V(d)=B^2\Var(\rho^4)+d^2\Var(\rho^2)
+2Bd\,\Cov(\rho^4,\rho^2).
\]
Using disk averages,
\[
\Var(\rho^4)=\frac4{45},\quad
\Var(\rho^2)=\frac1{12},\quad
\Cov(\rho^4,\rho^2)=\frac14-\frac16=\frac1{12}.
\]
Differentiate:
$V'(d)=d/6+B/6=0$, giving $d=-B$.
After piston removal the balanced residual is
\[
W_{\rm bal}=B(\rho^4-\rho^2+1/6)
=\frac{B}{6\sqrt5}Z_4^0,
\]
so
\begin{equation}
\sigma_{\rm bal}=\frac{|B|}{\sqrt{180}}.
\end{equation}
The original piston-removed quartic had
$\sigma=2|B|/\sqrt{45}$, exactly four times larger.
Allowing refocus changes the budget. That change is meaningful only
if refocus is an available adjustment under the actual imaging requirement.

\section{Sampling, obscurations, and amplitude weighting}
Real measured pupils contain discrete points and may have missing regions.
For samples $\vct w$, let $B_{\ell j}=Z_j(\rho_\ell,\theta_\ell)$ and
let $Q$ be a diagonal matrix of area or uncertainty weights.
Weighted least squares minimizes
$(B\vct c-\vct w)^TQ(B\vct c-\vct w)$ and gives
\begin{equation}
(B^TQB)\vct c=B^TQ\vct w.
\label{eq:modalfit}
\end{equation}
For a masked or nonuniformly weighted pupil, $B^TQB$ is generally not
the identity. Modes orthogonal on a full disk need not remain orthogonal
under that new metric.

Pupil \emph{intensity} weighting in wavefront metrology and pupil
\emph{amplitude} weighting in coherent on-axis interference are also
different. State which quantity is weighted before interpreting RMS,
Strehl, or fitted coefficients.

\checkpoint{Why can you not square and sum an arbitrary exported list
of coefficients to get RMS? You must first know basis normalization,
pupil support and weight, units, and which modes were removed.}
